\documentclass[11pt]{article}
\usepackage[T1]{fontenc}
\usepackage[a4paper,margin=24mm]{geometry}
\usepackage{amsmath,amssymb,amsthm}
\usepackage[hidelinks]{hyperref}
\usepackage{microtype}

\newtheorem{theorem}{Theorem}
\newcommand{\cH}{\mathcal H}
\newcommand{\cW}{\mathcal W}
\newcommand{\widehatGs}{\widehat G_{s}}

\setlength{\parindent}{0pt}
\setlength{\parskip}{0.56em}

\title{Wavelet Completeness Is Atomic Regularity}
\author{Literature-implied identification for arXiv:2005.04964}
\date{11 August 2026}

\begin{document}
\maketitle

\begin{center}
\fbox{\parbox{0.91\textwidth}{
\textbf{Status: literature-implied answer (exact conceptual
characterization).} The source's wavelet-complete locally compact groups are
exactly the established atomic-regular, or $[\mathrm{AR}]$, groups. Equivalently,
the Fourier algebra $A(G)$ is a dual Banach space, or $A(G)$ has the
Radon--Nikodym property. No new theorem is claimed.
}}
\end{center}

\section*{Source question}

In Section~7, arXiv PDF page~16, Berge asks to characterize the locally compact
groups $G$ satisfying
\begin{equation}\label{eq:source}
 \overline{\bigoplus_{\pi\in\widehatGs}
 \operatorname{span}_{g\in\mathcal A_\pi}
 \{\cW_gf:f\in\cH_\pi\}}=L^2(G),
\end{equation}
where $\widehatGs$ is the set of square-integrable irreducible representations,
$\mathcal A_\pi$ is the set of admissible vectors modulo phase, and
$\cW_gf(x)=\langle f,\pi(x)g\rangle$~\cite{berge}. The source calls such a group
\emph{wavelet complete}.

\section*{Identification}

Let $\lambda_G$ be the left regular representation on $L^2(G)$, and let
$L^2(G)_{\mathrm{at}}$ denote its atomic part: the closed sum of all irreducible
$\lambda_G$-invariant subspaces.

\begin{theorem}\label{thm:id}
For a locally compact group $G$,
\[
 \overline{\bigoplus_{\pi\in\widehatGs}
 \operatorname{span}_{g\in\mathcal A_\pi}\cW_g(\cH_\pi)}
 =L^2(G)_{\mathrm{at}}.
\]
Consequently the following conditions are equivalent:
\begin{enumerate}
\item $G$ is wavelet complete in the sense of \eqref{eq:source};
\item $\lambda_G$ is a Hilbert direct sum of irreducible representations;
\item $G$ is an atomic-regular ($[\mathrm{AR}]$) group.
\end{enumerate}
Under the standard hypotheses of the Fourier-algebra characterization, these
are also equivalent to $A(G)$ being a dual Banach space and to $A(G)$ having
the Radon--Nikodym property.
\end{theorem}

\emph{Why the theorem is already in the literature.}
The Duflo--Moore theorem gives, for every square-integrable irreducible $\pi$
and normalized admissible $g$, an isometric intertwiner
\[
 \cW_g:\cH_\pi\longrightarrow L^2(G),\qquad
 \cW_g\pi(y)=\lambda_G(y)\cW_g.
\]
Its range is therefore an irreducible subrepresentation of $\lambda_G$.
Conversely, every irreducible subrepresentation of $\lambda_G$ is
square-integrable; the Duflo--Moore discrete-part theorem identifies its
intertwining embeddings with coefficient transforms. Varying the admissible
vector supplies the multiplicity directions. Thus the closed coefficient sum
in \eqref{eq:source} is precisely $L^2(G)_{\mathrm{at}}$.

Taylor's survey states both directions on PDF page~5, then immediately defines
an $[\mathrm{AR}]$ group to be one whose left regular representation is a
direct sum of irreducibles~\cite{taylor}. On PDF page~8 it records the two
equivalent Fourier-algebra characterizations in Theorems~6 and~7. Hence
Theorem~\ref{thm:id} is a direct identification of the 2020 terminology with
a class studied since the 1970s, not a new result.

\section*{Consequences and scope}

The identification sharpens both examples in the source. Compact groups are
$[\mathrm{AR}]$ by Peter--Weyl, but compactness is not necessary: Taylor gives
the real affine group as a noncompact, nonunimodular $[\mathrm{AR}]$ example,
whose regular representation is a countable direct sum of two
square-integrable irreducibles. The reduced Heisenberg groups have a nonzero
continuous regular part, agreeing with the source's negative result.

For second-countable unimodular type-I groups, the same statement reads
particularly concretely:
\[
 G\text{ is wavelet complete}
 \quad\Longleftrightarrow\quad
 \text{Plancherel measure is purely atomic}.
\]
Each Plancherel atom is a discrete-series representation, and its matrix
coefficients span the corresponding Hilbert--Schmidt fiber.

This is an exact representation-theoretic and Banach-space characterization.
It is not a simple structural enumeration of every locally compact
$[\mathrm{AR}]$ group; that older classification program remains nontrivial.
The separate global HRT conjecture in the source is unaffected and remains
open.

\section*{Search and provenance}

The run indexes were searched for arXiv:2005.04964, HRT, wavelet completeness,
Plancherel atoms, and atomic regularity. Targeted searches for ``wavelet
complete'', ``completely reducible regular representation'', ``AR-group'', and
``purely atomic Plancherel measure'' located the older $[\mathrm{AR}]$
literature. Taylor's 2008 chapter predates arXiv:2005.04964 and therefore does
not explicitly claim to answer Berge's question. The implication is
agent-identified, so the correct provenance is \emph{literature implied}, not
\emph{literature already answered}.

\begin{thebibliography}{9}
\bibitem{berge}
E. Berge,
\emph{Interpolation in Wavelet Spaces and the HRT-Conjecture},
J. Fourier Anal. Appl. \textbf{26} (2020), article 72;
arXiv:2005.04964.

\bibitem{duflomoore}
M. Duflo and C. C. Moore,
\emph{On the regular representation of a nonunimodular locally compact
group}, J. Funct. Anal. \textbf{21} (1976), 209--243.

\bibitem{taylor}
K. F. Taylor,
\emph{Groups with Atomic Regular Representation}, in
\emph{Representations, Wavelets, and Frames}, Applied and Numerical Harmonic
Analysis, Birkh\"auser Boston (2008), 33--45;
doi:10.1007/978-0-8176-4683-7\_3.
\end{thebibliography}

\end{document}
